\section{A convergent joint crossover at order three}
\label{sec:order-crossover}

The fixed-order laws change at $s=3$, while the linear term in
\eqref{eq:normalform} changes at $b=\zeta(s-1)$. We now resolve both
transitions in one convergent chart. Throughout this section set
\begin{equation}\label{oc:paired-functions}
 A(\eps)=\zeta(2+\eps),\quad A_0=\zeta(2),\quad
 E_2(\eps)=\frac{\Gamma(1-\eps)}{(1+\eps)(1+\eps/2)},\quad
 Z(\eps)=\zeta(1+\eps)-\eps^{-1}.
\end{equation}
The last two functions are holomorphic near zero, with
$E_2(0)=1$, $E_2'(0)=\gamma-3/2$, and $Z(0)=\gamma$.
Pairing the resonant terms in \eqref{eq:jonquiere} gives
\begin{align}
 2K_{3+\eps}(x)&=x^2\{H(\eps,x)+xT(\eps,x)\},
       \label{oc:paired-K}\\
 H(\eps,x)&=1+\zeta(1+\eps)-\frac{E_2(\eps)x^\eps}{\eps},
       \label{oc:H}\\
 T(\eps,x)&=2\sum_{k\ge3}
       \frac{(-1)^k\zeta(3+\eps-k)}{k!}x^{k-3}.
       \label{oc:tail}
\end{align}
Here $T$ is jointly holomorphic for small $\eps$ and $|x|<2\pi$;
the paired Jonqui\`ere expansion converges normally on compact
subsets. In particular,
\begin{equation}\label{oc:paired-values}
 H(0,x)=\log(1/x)+\frac52,\qquad
 T(\eps,0)=-\frac{\zeta(\eps)}3,\qquad T(0,0)=\frac16.
\end{equation}
Thus \eqref{eq:normalform} becomes the exact equation
\begin{equation}\label{oc:exact-equation}
 d(2b+2x-d)=2\{b-A(\eps)\}x+x^2\{H(\eps,x)+xT(\eps,x)\}.
\end{equation}
The extra $1$ in \eqref{oc:H} records the quadratic feedback term;
it changes the resonant constant from $3/2$ to $5/2$.

\subsection{The joint order and feedback chart}

Put
\begin{equation}\label{oc:coordinates}
 X=\frac12\log\frac1d,\quad
 h=X^{-1},\quad w=X^{-1}\log X,\quad j=\sqrt{dX},\quad
 \eps=\lambda h,\quad b=A(\eps)+\eta j.
\end{equation}
In particular $d=hj^2$. The scaled displacement
$\eta=(b-A)/\sqrt{dX}$ measures when the linear and resonant
quadratic terms compete. Write
\begin{equation}\label{oc:profile}
 \ph(\lambda)=\frac{1-\ee^{-\lambda}}\lambda,\quad \ph(0)=1,
 \qquad
 S=\sqrt{\eta^2+2A_0\ph(\lambda)},\quad
 V_0=\frac{2A_0}{S+\eta}.
\end{equation}
For real $\lambda,\eta$, these quantities are positive. Define
\begin{equation}\label{oc:G}
 G(\lambda,h,w,V)=
 \frac{1-\ee^{-\lambda}E_2(\lambda h)
           \ee^{-\lambda w/2}V^{\lambda h}}\lambda
       +h\{1+Z(\lambda h)\},
\end{equation}
using the logarithm near $V_0>0$. Its quotient is holomorphic at
$\lambda=0$, where
\begin{equation}\label{oc:G-zero}
 G(0,h,w,V)=1+\frac w2+h\left(\frac52-\log V\right).
\end{equation}

\begin{theorem}[Convergent joint crossover]\label{thm:joint-crossover}
For every compact rectangle of real parameters $(\lambda,\eta)$,
there is a function $V(\lambda,\eta,h,w,j)$ holomorphic on a common
complex neighborhood of that rectangle times a polydisc about
$(h,w,j)=(0,0,0)$, with $V(\lambda,\eta,0,0,0)=V_0$.
For all sufficiently large real $X$, uniformly on the rectangle,
the positive feedback inverse in \eqref{oc:coordinates} is exactly
\begin{equation}\label{oc:joint-chart}
 \boxed{\displaystyle
 x=\sqrt{\frac dX}\,
 V\left(\lambda,\eta,\frac1X,\frac{\log X}{X},\sqrt{dX}\right).}
\end{equation}
Its defining holomorphic equation is
\begin{align}
 0={}&V^2G(\lambda,h,w,V)+2\eta V
      +h^2jV^3T(\lambda h,hjV)\notag\\
    &-2A(\lambda h)-2j\eta-2hjV+hj^2.
       \label{oc:joint-implicit}
\end{align}
Thus the Taylor series in the three independent small coordinates
converges after physical substitution; it includes every
inverse-logarithmic order and all exponential sectors.
\end{theorem}

\begin{proof}
At $h=w=j=0$, the right side of \eqref{oc:joint-implicit} is
$\ph(\lambda)V^2+2\eta V-2A_0$. Its positive root is $V_0$, and
its derivative there is $2S\ne0$. Joint holomorphy, the implicit
function theorem, and compactness give a common polydisc; uniqueness
glues the local functions. Under $x=hjV$, the physical coordinates
satisfy
\[
 x^\eps=\ee^{-\lambda}\ee^{-\lambda w/2}V^{\lambda h},
 \qquad hH(\eps,x)=G(\lambda,h,w,V).
\]
Dividing \eqref{oc:exact-equation} by $d=hj^2$ now gives
\eqref{oc:joint-implicit} exactly. For the constructed root,
$b+x-d$ is close to $A_0>0$, so the required square-root branch in
\eqref{eq:Psi} is preserved.

For completeness the positive-root identification remains valid
when $\eta<0$. For real $s>2$ and $b>0$, the ratio
\[
 \frac{\Li_{s-1}(\ee^{-x})}{\sqrt{b^2+2D_s(x)}}
\]
strictly decreases from $A/b$ to zero as $x$ increases from zero
to infinity. Thus $\Psi_b'(x)$ strictly increases to $1$.
If $b\ge A$, then $\Psi_b$ increases from zero. If $b<A$, it
first decreases to a negative minimum and then increases to
infinity. In either case every $d>0$ has exactly one positive
preimage. In the latter case this root lies before the interior
fold when the branch is followed from $r=0$, since $r<\Rb$;
the endpoint reference $\Rb$ need not itself be the Taylor radius.
The positive root furnished by the chart is therefore the physical
one. Finally, the asserted convergence is the ordinary convergence
of the holomorphic Taylor series, not an inference from its
formal substitutions.
\end{proof}

\begin{corollary}[Explicit first correction]\label{cor:joint-first}
Uniformly on compact real $(\lambda,\eta)$ sets, let
\begin{equation}\label{oc:first-B}
 B=\ee^{-\lambda}
   \left\{\frac12\log X-\log V_0-\gamma+\frac32\right\}
      +1+\gamma,
 \qquad V_1=\frac{2\lambda\zeta'(2)-V_0^2B}{2S}.
\end{equation}
Then
\begin{equation}\label{oc:first-correction}
 x=\sqrt{\frac dX}
   \left\{V_0+\frac{V_1}{X}
       +O\left(\frac{(1+\log X)^2}{X^2}\right)\right\}.
\end{equation}
More generally, the $j=0$ part of \eqref{oc:joint-chart} is a
convergent series $\sum_{n\ge0}X^{-n}P_n(\log X;\lambda,\eta)$
with $\deg P_n\le n$. Truncation after $n=N$ has uniform error
$O((1+\log X)^{N+1}X^{-N-1})$ inside the braces.
\end{corollary}

\begin{proof}
Expand $\log E_2(\eps)=(\gamma-3/2)\eps+O(\eps^2)$ and
$A(\eps)=A_0+\zeta'(2)\eps+O(\eps^2)$ in
\eqref{oc:joint-implicit}, first setting $j=0$. Its coefficient
of $X^{-1}$ is $2SV_1+V_0^2B-2\lambda\zeta'(2)$.
The remaining assertions follow by grouping the absolutely
convergent Taylor series in $(h,w)$ by total degree. A monomial
$h^i w^j$ contributes $X^{-i-j}(\log X)^j$, which proves both
the degree bound and the Taylor remainder. Terms containing the
third coordinate $\sqrt{dX}$ are smaller than every fixed inverse
power of $X$ and are retained in the full chart.
\end{proof}

At $\lambda=\eta=0$, this recovers the explicit resonance constant:
\begin{equation}\label{oc:critical-first}
 x=\sqrt{\frac{2\zeta(2)d}{X}}
 \left\{1-\frac{\log X-\log(2\zeta(2))+5}{4X}
       +O\left(\frac{(1+\log X)^2}{X^2}\right)\right\}.
\end{equation}
The compact-parameter hypothesis is essential to the stated
uniform errors; it does not by itself permit $\lambda$ or $\eta$
to tend to infinity with $X$.

\subsection{Critical support and the first exponential correction}

At $\eta=0$, different analytic coordinates reveal an additional
restriction on the support. Set $q=X^{-1/2}$ and $z=\sqrt d$, so
$x=zqU$. There is a holomorphic chart $U(\lambda,q,w,z)$ through
$U_0=\sqrt{2A_0/\ph(\lambda)}$, determined by
\begin{align}
 0={}&U^2G(\lambda,q^2,w,U)
       +zq^3U^3T(\lambda q^2,zqU)\notag\\
    &-2A(\lambda q^2)-2zqU+z^2.
       \label{oc:critical-implicit}
\end{align}
Its existence follows directly from the same implicit-function
argument, since the derivative at the base root is
$2\ph(\lambda)U_0\ne0$.

\begin{proposition}[Parity and the first exponential sector]
\label{prop:critical-support}
The critical chart satisfies
\begin{equation}\label{oc:parity}
 U(\lambda,-q,w,-z)=U(\lambda,q,w,z).
\end{equation}
Hence the coefficient of $q^iw^jz^k$ vanishes when $i+k$ is odd.
Put $h=q^2$ and, at $z=0$, write
\[
 U_*=U(\lambda,q,w,0),\quad G_*=G(\lambda,h,w,U_*),\quad
 M_*=\ee^{-\lambda}E_2(\lambda h)
             \ee^{-\lambda w/2}U_*^{\lambda h}.
\]
The exact first exponential coefficient is
\begin{equation}\label{oc:first-exp-coefficient}
 \partial_zU(\lambda,q,w,0)
 =q\,\frac{2+hU_*^2\zeta(\lambda h)/3}{2G_*-hM_*}.
\end{equation}
Consequently, for $x_*=\sqrt{d/X}\,U_*$ on the physical curve,
\begin{equation}\label{oc:first-exp-sector}
 x=x_*+\frac dX
       \frac{2+X^{-1}U_*^2\zeta(\lambda/X)/3}
            {2G_*-X^{-1}M_*}
       +O\left(d^{3/2}X^{-1/2}\right).
\end{equation}
The correction begins with $d/(X\ph(\lambda))$. All statements
are uniform for $\lambda$ in compact real sets.
\end{proposition}

\begin{proof}
The defining equation \eqref{oc:critical-implicit} is invariant
under $(q,z)\mapsto(-q,-z)$; uniqueness proves \eqref{oc:parity}.
At $z=0$, its $U$ derivative equals
$U_*(2G_*-hM_*)$, because $G_U=-hM_*/U_*$ there.
Its $z$ derivative equals
$-q^3U_*^3\zeta(\lambda h)/3-2qU_*$, by
\eqref{oc:paired-values}. Implicit differentiation proves
\eqref{oc:first-exp-coefficient}. Taylor's theorem in $z$, followed
by multiplication by $zq$, gives \eqref{oc:first-exp-sector}.
The parity statement also independently explains why the
$z=0$ sector contains integer powers of $X^{-1}$.
\end{proof}

\subsection{A local complex fold beyond a profile zero}

The real crossover profile is positive, but its nonzero complex
zeros are $\lambda_0=2\pi\ii k$, $k\ne0$. The following theorem
concerns local solutions of the continued normal form at $b=A$.
It makes no assertion that those solutions continue the Taylor germ
$v(0)=0$ around any specified global path.

\begin{theorem}[Local shifted complex fold]\label{thm:local-complex-fold}
Fix $\lambda_0=2\pi\ii k\ne0$, and put
$h=1/X$, $z=\sqrt d$, $\lambda=\lambda_0+\alpha h$,
$\eps=h(\lambda_0+\alpha h)$, and $x=z\omega$.
There are holomorphic functions $\alpha_c(h,z)$ and $\omega_c(h,z)$
near $(0,0)$ with
\begin{equation}\label{oc:complex-base}
 \omega_c(0,0)=2\sqrt{A_0},\qquad
 \alpha_c(0,0)=\alpha_*:=
 \lambda_0\left\{\frac12\log(4A_0)-2\right\},
\end{equation}
on which the local inverse equation has a double root.
An analytic local coordinate converts that equation into a square
equal to $(\alpha-\alpha_c(h,z))$ times a nonvanishing holomorphic
factor. In particular, the scaled-order fold curve satisfies
\begin{equation}\label{oc:complex-curve}
 \lambda_c(X)=\lambda_0+
 \frac{\lambda_0\{\frac12\log(4A_0)-2\}}{X}
       +O(X^{-2}+X^{-1}\ee^{-X}).
\end{equation}
At the unshifted profile zero $\alpha=0$, there are instead two
ordinary holomorphic root germs in $(h,z)$, through the distinct
positive constants
\begin{equation}\label{oc:complex-unshifted}
 \omega_j=\exp\left\{\frac{5+W_j(-4A_0\ee^{-5})}{2}\right\},
 \qquad j=0,-1.
\end{equation}
Thus the unshifted profile zero itself does not force a
square-root inverse-logarithmic splitting.
\end{theorem}

\begin{proof}
Define, with the logarithm near the indicated positive values,
\[
 \widehat H(h,\alpha,\omega)=1+Z(\eps)
    +\frac{1-E_2(\eps)\ee^{-\alpha h}\omega^\eps}{\eps}.
\]
The numerator is divisible by $h$, so this expression is
holomorphic at $h=0$. Its value there is
$5/2+\alpha/\lambda_0-\log\omega$. The exact equation is
\begin{equation}\label{oc:complex-equation}
 \widehat F:=\omega^2\widehat H+z\omega^3T(\eps,z\omega)
               -2A(\eps)-2z\omega+z^2=0.
\end{equation}
At $h=z=0$ it reduces to
\[
 \omega^2\left(\frac52+\frac\alpha{\lambda_0}-\log\omega\right)=2A_0.
\]
Its double positive root is exactly \eqref{oc:complex-base}; there
$\widehat F_{\omega\omega}=-2$ and
$\widehat F_\alpha=4A_0/\lambda_0\ne0$. Apply the implicit
function theorem to $(\widehat F,\widehat F_\omega)$ to obtain
$\alpha_c,\omega_c$. For the normal form, first solve
$\widehat F_\omega=0$ for a stationary point depending on
$(h,z,\alpha)$. Taylor expansion about it factors the remaining
quadratic term by a holomorphic unit. Taking its holomorphic square
root gives a square coordinate. The value at the stationary point
has a simple zero at $\alpha_c$, because its $\alpha$ derivative
is nonzero. This proves the claimed factorization and
\eqref{oc:complex-curve}.

At $\alpha=0$, the limiting equation is
$\omega^2(5/2-\log\omega)=2A_0$. Solving it gives
\eqref{oc:complex-unshifted}. Since $A_0<2$ and $8<\ee^4$,
its Lambert argument lies strictly between $-1/\ee$ and zero.
The two stated roots are therefore distinct and simple, and the
ordinary implicit function theorem supplies their local charts.
No exhaustion of other complex roots is claimed.
\end{proof}
